
Here, we provide the technical details of generating \cref{fig:ridge_regression_prob_dist_loss_landscape} in the main text, which exmplifies how the probability distribution of the hyperparameters changes at different points on the landscape of the loss function for model training, following the underlying example given in \citet{prince2023understanding}.

The simple ridge regression model is applied to fit the data generated from the \emph{Gabor model} function that maps an 1D input to an 1D output, which is defined as follows:
\begin{equation}
    \label{eq:gabor_model}
 y=f(x|\phi_0,\phi_1)=\sin\left(\phi_0+0.06\phi x\right)\cdot\exp\left( -\frac{(\phi_0 + 0.06\phi_1 x)^2}{32.0}\right),
\end{equation}
in which $\phi_0$ and $\phi_1$ are the hyperparameters of the Gabor model, that we shorthanded it with a vectorial quantity $\boldsymbol{\phi}=(\phi_0,\phi_1)^\top$. We generated 40 testing samples by first randomly sampled 40 $x$ values in the interval $[-15,15]$ and then computed the corresponding $y$ values using \cref{eq:gabor_model} with additional noises $\epsilon\sim\mathrm{Uniform}(-0.2,0.2)$ added.

The landscape of the loss function is generated by computing the values of the loss function over a uniform 2D grid of $(\phi_0,\phi_1)$. We take the loss function for the ridge regression:
\begin{equation*}
    \label{eq:gabor_ridge_loss}
    \mathcal{L}(\boldsymbol{\phi})=\sum_{i=1}^{N}\left\Vert y_i-f(x_i|\boldsymbol{\phi} )\right\Vert_2^2+\lambda\left\Vert \boldsymbol{\phi} \right\Vert_2^2,
\end{equation*}
in which we set the regularisation strength $\lambda=10^{-4}$. 

Unfortunately, computing $\mathcal{L}(\boldsymbol{\phi})$ over the $\boldsymbol{\phi}$-grid does not give us the direct access to the probability distribution of the hyperparameters $p(\boldsymbol{\phi})$, which is what we want to show in \cref{fig:ridge_regression_prob_dist_loss_landscape}. This problem is overcomed by realising the fact that the local curvature of the loss function at a given point $\hat{\boldsymbol{\phi}}$ can be used to obtain a Boltzmann-like distribution of the hyperparameters around $\hat{\boldsymbol{\phi}}$ as $p(\hat{\boldsymbol{\phi}})\propto \exp(-\beta \hat{\boldsymbol{\phi}})$, which is approximately a Gaussian distribution $\mathcal{N}(\hat{\boldsymbol{\phi}},\boldsymbol{\Sigma})$ with mean $\hat{\boldsymbol{\phi}}$ and a covariance matrix $\boldsymbol{\Sigma}\sim(\beta\bm{H})^{-1}$, if we perform a second-order Taylor expansion of the loss function around $\hat{\boldsymbol{\phi}}$:
\begin{equation*}
    \label{eq:taylor_expansion_loss}
    \mathcal{L}(\boldsymbol{\phi})\approx \mathcal{L}(\hat{\boldsymbol{\phi}})+\frac{1}{2}(\boldsymbol{\phi}-\hat{\boldsymbol{\phi}})^\top \bm{H} (\boldsymbol{\phi}-\hat{\boldsymbol{\phi}}),
\end{equation*}
in which $\bm{H}$ is the Hessian matrix of the loss function at $\hat{\boldsymbol{\phi}}$, which can be computed using a finite difference method around the vicinity of $\hat{\boldsymbol{\phi}}$. 

Nevertheless, for an arbitrary point $\hat{\boldsymbol{\phi}}$ on the loss landscape, the Hessian matrix $\bm{H}$ may not be positively definitive and hence cannot be directly inverted. To solve this problem, we perform an eigen-decomposition of $\bm{H}$ and then set the negative eigenvalues to a small positive value of $10^{-3}$ to recompute the modified Hessian matrix $\tilde{\bm{H}}$ to make it become positively definitive, from which we can compute the covariance matrix $\boldsymbol{\Sigma}\sim(\beta\tilde{\bm{H}})^{-1}$ and visualise the probability distribution $\mathcal{N}(\hat{\boldsymbol{\phi}},\boldsymbol{\Sigma})$ as shown in \cref{fig:ridge_regression_prob_dist_loss_landscape}. Here, the temperature $\beta$ is set to $10^{3}$ to make the probability distribution more concentrated around $\hat{\boldsymbol{\phi}}$ for better visualisation.